%--------------------------------------------------------------------------
% Angela R. Lopez Sanchez - Hoja de Vida (Espanol)
% Compilar con: xelatex CV_Lopez_ES.tex
% Formato adaptado de Camacho_CV_Base.tex
%--------------------------------------------------------------------------
\documentclass[10pt]{article}

%----- Packages ------------------------------------------------------------
\usepackage[letterpaper,
            top=0.52in, bottom=0.52in,
            left=0.52in, right=0.52in]{geometry}

\usepackage{fontspec}           % xelatex font selection
\usepackage{microtype}
\usepackage{parskip}
\usepackage{enumitem}
\usepackage{titlesec}
\usepackage{xcolor}
\usepackage[spanish,es-noindentfirst]{babel}
\usepackage[hidelinks]{hyperref}

%----- Font ----------------------------------------------------------------
\setmainfont{Times New Roman}

%----- Colors & Hyperlinks -------------------------------------------------
\definecolor{urlblue}{RGB}{5,99,193} % Standard Word hyperlink blue
\hypersetup{
  colorlinks=true,
  urlcolor=urlblue,
  pdfauthor={Angela R. Lopez Sanchez},
  pdftitle={Angela R. Lopez Sanchez - Hoja de Vida}
}

%----- Global spacing ------------------------------------------------------
\setlength{\parskip}{0pt}
\setlength{\parindent}{0pt}
\pagestyle{empty}

%----- Section heading style -----------------------------------------------
\titleformat{\section}
  {\normalfont\large\bfseries\MakeUppercase}
  {}{0em}{}
  [\vspace{-10pt}\rule{\linewidth}{0.5pt}\vspace{-2pt}]
\titlespacing*{\section}{0pt}{10pt}{0pt}

%----- Reusable layout commands --------------------------------------------
\newcommand{\org}[2]{\noindent\textbf{#1}\hfill #2\par\vspace{0pt}}
\newcommand{\role}[2]{\noindent\textit{#1}\hfill #2\par\vspace{0pt}}
\newcommand{\edu}[2]{\noindent #1\hfill #2\par\vspace{0pt}}

% Bullet list style
\newenvironment{cvitems}{%
  \begin{itemize}[
    leftmargin=1.2em,
    itemsep=0pt,
    parsep=0pt,
    topsep=1pt,
    label={$\bullet$}]%
}{\end{itemize}}

\newcommand{\entrygap}{\vspace{6pt}}

%==========================================================================
\begin{document}
%==========================================================================

%------ ENCABEZADO ---------------------------------------------------------
\begin{center}
  {\LARGE\textbf{ANGELA R. LOPEZ SANCHEZ}}\\[4pt]
  +1 (202) 390-4849 $|$
  \href{mailto:alopezsanchezv@gmail.com}{alopezsanchezv@gmail.com} $|$
  \href{https://github.com/AngelaLop}{GitHub} $|$
  \href{https://angelalop.github.io/AngelaLopezS/}{Portafolio} $|$
  \href{https://scholar.google.com/citations?user=wfO9incAAAAJ&hl=en}{Google Scholar}
\end{center}

\vspace{3pt}

%------ EDUCACION ----------------------------------------------------------
\section{Educaci\'on}

\org{The University of Chicago}{Chicago, EE. UU.}
\edu{Maestr\'ia en An\'alisis Computacional y Pol\'iticas P\'ublicas}{Junio de 2026}
\textbf{Cursos relevantes:} Aprendizaje Autom\'atico Avanzado para Pol\'iticas P\'ublicas, Ciencias de la Computaci\'on con Aplicaciones, Computaci\'on en la Nube, Bases de Datos, Visualizaci\'on de Datos\\
\textbf{Reconocimiento:} Banco de la Rep\'ublica, Beca Nacional al M\'erito, una de solo dos otorgadas a nivel nacional para estudios de maestr\'ia en pol\'iticas p\'ublicas

\entrygap
\org{Universidad de los Andes}{Bogot\'a, Colombia}
\edu{Maestr\'ia en Ciencias Sociales con \'enfasis en Estudios Internacionales}{Abril de 2018}
\textit{Tesis:} Evaluaci\'on de impacto del retorno de colombianos desde Venezuela sobre el bienestar de los hogares receptores

\entrygap
\org{Universidad Militar Nueva Granada}{Bogot\'a, Colombia}
\edu{Pregrado en Econom\'ia (Economista)}{Octubre de 2011}

%------ HABILIDADES --------------------------------------------------------
\section{Habilidades}

\textbf{Lenguajes de Programaci\'on:} Python, R, Stata, SQL, JavaScript, HTML.
\textbf{Librer\'ias y Frameworks:} D3.js, Altair/Vega, Dash, Next.js, MapLibre.
\textbf{Nube e Infraestructura de Datos:} AWS (S3, EC2, IAM), DuckDB, Supabase, Git/GitHub, pipelines reproducibles en CLI.
\textbf{Geoespacial y Percepci\'on Remota:} Google Earth Engine, QGIS, ArcGIS, OSRM, FMM.
\textbf{M\'etodos:} Aprendizaje autom\'atico, inferencia causal y econometr\'ia, an\'alisis geoespacial, microsimulaci\'on, datos de encuestas.
\textbf{BI:} Power BI.\\
\textbf{Idiomas:} Espa\~nol (Nativo), Ingl\'es (Fluido), Franc\'es (B\'asico)

%------ EXPERIENCIA PROFESIONAL --------------------------------------------
\section{Experiencia Profesional}

\org{Banco Interamericano de Desarrollo (BID)}{Bogot\'a, Colombia (Remoto)}
\role{Consultora Senior en An\'alisis Geoespacial}{Mayo de 2025 -- Presente}
\vspace{-1pt}
\begin{cvitems}
  \item Dise\~n\'e y entregu\'e el primer conjunto de datos regional de accesibilidad escolar para Am\'erica Latina y el Caribe (22 pa\'ises, 530.000 escuelas), hoy utilizado por los equipos de educaci\'on del BID para medir brechas de acceso
  \item Constru\'i un pipeline reproducible en Python (CLI) que abarca geocodificaci\'on, control de calidad y modelaci\'on de tiempos de viaje por dos m\'etodos (FMM y enrutamiento de red OSRM) para calcular la accesibilidad a escala continental
\end{cvitems}

\entrygap
\org{Organizaci\'on para la Cooperaci\'on y el Desarrollo Econ\'omicos (OCDE)}{Par\'is, Francia}
\role{Consultora en Datos y Pol\'iticas}{Junio -- Septiembre de 2025}
\vspace{-1pt}
\begin{cvitems}
  \item Consolid\'e el primer conjunto de datos de capital social de la OCDE a nivel subnacional (TL2/TL3) para los pa\'ises miembros
  \item Realic\'e el an\'alisis estad\'istico y dise\~n\'e todas las visualizaciones de datos del Cap\'itulo 3, ``Desigualdades geogr\'aficas en el acceso a oportunidades'', del informe insignia \textit{To Have and Have Not: How to Bridge the Gap in Opportunities}
\end{cvitems}

\entrygap
\org{Banco Mundial}{Washington, D.C.}
\role{Consultora en Pobreza}{Agosto de 2021 -- Agosto de 2024}
\vspace{-1pt}
\begin{cvitems}
  \item Simul\'e resultados de pobreza y desigualdad en Panam\'a bajo distintos escenarios de pol\'itica y macroecon\'omicos usando modelos de microsimulaci\'on en Stata, como insumo para las Reuniones Anuales y de Primavera del Banco
  \item Evalu\'e el impacto distributivo de programas sociales (p.\,ej., transferencias monetarias) bajo distintos criterios de elegibilidad para orientar la focalizaci\'on de la poblaci\'on beneficiaria en Panam\'a
  \item Elabor\'e y contribu\'i a 3 productos anal\'iticos sobre pobreza, desigualdad, brechas de g\'enero y exclusi\'on para apoyar el di\'alogo de pol\'itica con el gobierno, incluyendo una actualizaci\'on del Diagn\'ostico Sistem\'atico de Pa\'is y una Evaluaci\'on de Pobreza en la que lider\'e el cap\'itulo de fondo sobre mercado laboral
\end{cvitems}

\entrygap
\org{Banco Interamericano de Desarrollo (BID)}{Washington, D.C.}
\role{Consultora en Datos y Pol\'iticas}{Junio -- Noviembre de 2022}
\vspace{-1pt}
\begin{cvitems}
  \item Consolid\'e el primer conjunto de datos regional sobre las caracter\'isticas de las evaluaciones de aprendizaje para Am\'erica Latina y el Caribe
  \item Redact\'e y entregu\'e el informe de pol\'itica \textit{El estado de la educaci\'on en Am\'erica Latina y el Caribe: las evaluaciones de aprendizaje}, como insumo para el di\'alogo de pol\'itica educativa basado en evidencia
\end{cvitems}

\entrygap
\org{Banco Interamericano de Desarrollo (BID)}{Washington, D.C.}
\role{Analista Econ\'omica, Divisi\'on de Educaci\'on}{Julio de 2018 -- Agosto de 2021}
\vspace{-1pt}
\begin{cvitems}
  \item Produje estad\'isticas de educaci\'on y sociales para pa\'ises de Am\'erica Latina y el Caribe a partir de encuestas de hogares y datos administrativos, generando productos anal\'iticos que apoyaron el di\'alogo de pol\'itica con gobiernos de la regi\'on
  \item Brind\'e apoyo anal\'itico y estad\'istico para poner en marcha la nueva agenda de Educaci\'on Superior del Banco para Am\'erica Latina
  \item Desarroll\'e herramientas de datos (un tablero y reportes automatizados) y asesor\'e a las Divisiones del Sector Social del Banco en an\'alisis de datos recurrente
\end{cvitems}

\entrygap
\org{Departamento Administrativo Nacional de Estad\'istica (DANE)}{Bogot\'a, Colombia}
\role{Coordinadora, Unidad de ODS}{Febrero -- Junio de 2018}
\vspace{-1pt}
\begin{cvitems}
  \item Lider\'e un equipo de 4 personas en distintas \'areas del DANE para definir y dise\~nar la metodolog\'ia de 56 indicadores nacionales de seguimiento a los Objetivos de Desarrollo Sostenible
  \item Coordin\'e con entidades externas para cumplir los requerimientos de reporte de la Agenda Post-2015 y represent\'e a la oficina de estad\'istica de Colombia en 3 reuniones internacionales sobre metodolog\'ias de indicadores de los ODS
\end{cvitems}

\entrygap
\org{Departamento Administrativo Nacional de Estad\'istica (DANE)}{Bogot\'a, Colombia}
\role{Analista Econ\'omica, Unidad de Mercado Laboral}{Febrero de 2015 -- Enero de 2018}
\vspace{-1pt}
\begin{cvitems}
  \item Realic\'e an\'alisis mensuales del mercado laboral y de indicadores sociodemogr\'aficos de Colombia para difusi\'on nacional e internacional
  \item Constru\'i y proces\'e las bases mensuales de la Encuesta de Hogares que sustentan las estad\'isticas oficiales del mercado laboral y socioecon\'omicas, y capacit\'e trimestralmente a encuestadores para asegurar la calidad de los datos
\end{cvitems}

\entrygap
\org{Universidad Militar Nueva Granada}{Bogot\'a, Colombia}
\role{Asistente de Investigaci\'on}{Febrero de 2013 -- Diciembre de 2014}
\vspace{-1pt}
\begin{cvitems}
  \item Coautora de 7 ponencias acad\'emicas nacionales y 3 internacionales, 3 art\'iculos publicados y 8 objetos virtuales de aprendizaje
  \item Dirig\'i las sesiones semanales del Grupo de Estudios Econ\'omicos y asesor\'e a estudiantes de pregrado en proyectos de iniciaci\'on cient\'ifica
\end{cvitems}

%------ PROYECTOS DESTACADOS -----------------------------------------------
\section{Proyectos Destacados}

\begin{cvitems}
  \item \textbf{EduAccess LAC, Plataforma de Accesibilidad Escolar:} Desarroll\'e y publiqu\'e una geo-plataforma full-stack (Next.js, Supabase Postgres, worker cron en Railway) que convierte los indicadores de accesibilidad del BID para cinco pa\'ises de Am\'erica Latina y el Caribe en una herramienta que un director de ministerio sin perfil t\'ecnico puede usar; constru\'i una cascada de dos niveles de LLM (Llama 3.1/3.3) que traduce preguntas en lenguaje natural a SQL validado y protegido contra inyecci\'on de prompts, y resalta los resultados en un mapa MapLibre \href{https://eduaccess-lac.vercel.app/}{(En vivo)}
  \item \textbf{Clasificaci\'on de Pobreza con Datos Multimodales:} Entren\'e un modelo de aprendizaje autom\'atico que clasifica hogares en categor\'ias de pobreza a partir de im\'agenes satelitales y unas pocas preguntas de prueba de medios (proxy means test), una forma de menor costo para focalizar programas sociales donde faltan datos de encuestas (Python, ML) \href{https://github.com/uchicago-capp30254-spr-25/project-alopezsanchez-cnunezh-jmcarias}{(Repositorio)}
  \item \textbf{Explorador del Ciclo de Vida en Panam\'a:} Visualizaci\'on interactiva de c\'omo cambia la participaci\'on en educaci\'on, trabajo y jubilaci\'on seg\'un edad, g\'enero y regi\'on en Panam\'a, mostrando d\'onde los j\'ovenes abandonan la escuela o las mujeres salen del mercado laboral (Python, D3.js) \href{https://angelalop.github.io/capp30239_interactive_viz/}{(En vivo)}
\end{cvitems}

\vspace{2pt}
Para la lista completa de proyectos, visite mi \href{https://angelalop.github.io/AngelaLopezS/}{portafolio}.

%------ PUBLICACIONES ------------------------------------------------------
\section{Publicaciones Seleccionadas}

\textbf{Art\'iculos en Revistas Arbitradas}
\begin{cvitems}
  \item L\'opez, A.R., Okoye, K., Haruna, H., et al. (2022). Impact of Digital Technologies upon Teaching and Learning in Higher Education in Latin America: an Outlook on the Reach, Barriers, and Bottlenecks. \textit{Education and Information Technologies}
  \item L\'opez, A.R., Virg\"uez, A.F., Silva, A.C., \& Sarmiento, J.A. (2017). Desigualdad de oportunidades en el sistema de educaci\'on p\'ublica en Bogot\'a, Colombia. \textit{Lecturas de Econom\'ia}, (87), 165--190
  \item L\'opez, A.R., Virg\"uez, A.F., Sarmiento Espinel, J.A., \& Silva Arias, A.C. (2015). El efecto de la gerencia privada de escuelas p\'ublicas en el desempe\~no estudiantil en la educaci\'on media en Colombia. \textit{Ecos de Econom\'ia}, 19(41), 108--136
\end{cvitems}

\vspace{3pt}
\textbf{Notas de Pol\'itica (Selecci\'on)}
\begin{cvitems}
  \item L\'opez, A. (2024). \textit{Panama Poverty Assessment: Technical Note on Labor Market}. The World Bank, Washington, D.C.
  \item Arias, E., L\'opez, A., Due\~nas, X., \& Giambruno, C. (2024). \textit{El estado de la educaci\'on en Am\'erica Latina y el Caribe: las evaluaciones de aprendizaje}. Banco Interamericano de Desarrollo, Washington, D.C.
  \item Arias, E., L\'opez, A., Elacqua, G., et al. (2021). \textit{Educaci\'on superior y COVID-19 en Am\'erica Latina y el Caribe: Financiamiento para los estudiantes}. Banco Interamericano de Desarrollo, Washington, D.C.
\end{cvitems}

\vspace{2pt}
Para la lista completa de publicaciones, consulte mi \href{https://scholar.google.com/citations?user=wfO9incAAAAJ&hl=en}{perfil de Google Scholar}.

%==========================================================================
\end{document}
